\chapter{Evaluation Protocol}
\thispagestyle{empty}

The proposed pipeline and every method of the preliminary study (Chapters~5--7) are evaluated in
exactly the same way, described in this chapter. The
protocol is designed to measure what a detection system needs: how well it classifies a
person whose EEG it has never seen.

\section{Subject-wise cross-validation}
\label{sec:protocol}
All windows of a person are kept in the same fold, so the people in a test fold are unknown to
the model trained on the other folds. For DS-1, leave-one-subject-out (LOSO) cross-validation
is used: 14 folds, each with one participant as the test set and the other 13 for training.
For DS-2, stratified grouped 10-fold cross-validation is used: the 121 children are divided
into 10 folds with the classes balanced in each, using a fixed random seed, and each fold
serves once as the test set. After all folds, every person has been tested exactly once by a
model that never saw them, and all methods use the same folds, so their results can be
compared directly.

\section{No information from the test subjects}
Every step that learns from data is fitted on the training subjects of the fold only: feature
standardisation, feature selection, the reference point of the tangent space, the selection of
channels, the choice of hyperparameters and the early stopping of the networks. The only
exceptions are three methods of the preliminary study that by design use unlabelled windows of
the test subject; they are marked $^\dagger$ in its tables (Appendix~C) and are not used in the
proposed pipeline.

Hyperparameters are chosen by \emph{nested} cross-validation. Inside each outer training set,
an inner stratified grouped 4-fold cross-validation, again by subject, evaluates every point of
a small grid; the point with the highest pooled balanced accuracy is used to train on the whole
outer training set. The outer test subjects take no part in this choice.

An automated test suite checks these properties. It asserts that no subject appears in both the
training and the test part of any outer or inner fold, and it runs the nested procedure with a
classifier that records the subjects it is given, to confirm that no fit ever receives a window
of a test subject.

\section{Subject-level decisions}
A person is assigned the class predicted for the majority of their windows. Ties are broken by
the mean predicted probability and counted. Results are reported at both levels: the window
level shows how reliable a single 5~s decision is, and the subject level is the decision that
would be given for a person.

\section{Performance measures}
With TP, TN, FP and FN the numbers of true positives, true negatives, false positives and false
negatives, where the disorder is the positive class:
\begin{align}
\text{accuracy} &= \frac{TP+TN}{TP+TN+FP+FN}, &
\text{sensitivity} &= \frac{TP}{TP+FN}, \nonumber\\
\text{specificity} &= \frac{TN}{TN+FP}, &
\text{balanced accuracy} &= \frac{\text{sensitivity}+\text{specificity}}{2}, \\
\text{F1} &= \frac{2\,TP}{2\,TP+FP+FN}. \nonumber
\end{align}
Balanced accuracy is the main measure, because the classes of DS-2 have different numbers of
windows; a classifier that always predicts one class has a balanced accuracy of 50\%. The area
under the ROC curve (AUC) is the probability that a randomly chosen positive example receives a
higher score than a randomly chosen negative one; 0.5 is chance, and a value below 0.5 means
that the score is related to the class in the opposite direction.

\section{Confidence intervals and statistical checks}
\textbf{Bootstrap over subjects.} All intervals are 95\% percentile bootstrap
intervals~\cite{efron1993} over \emph{subjects}: subjects are resampled with replacement, within
each class, 2000 times, each resampled subject bringing all of their windows, and the measure is
recomputed for each resample. Resampling windows instead would treat the windows of one person
as independent and give intervals that are far too narrow.

If every subject is classified correctly, every resample is also entirely correct and the
bootstrap interval collapses to 100--100\%. In that case the exact binomial (Clopper--Pearson)
interval~\cite{clopper1934} is the appropriate statement of uncertainty; for 14 correct out of
14 it is \R{DSoneExactAllCorrectLow}--100\%.

\textbf{Selection among many variants.} When several variants are compared on the same test
folds, the best of them is slightly optimistic. Decisions that are part of the proposed pipeline,
such as which common modes to use, are therefore made by the inner cross-validation on the
training subjects, and the per-variant results of the ablations are reported as analyses, not as
the result of the pipeline.

\section{Why the windows of a person must not be split}
\label{sec:whysubject}
Many published results on these datasets, including the accuracies that made the SMVMD method
the starting point of this project, were obtained by shuffling all windows into folds
(\emph{window-wise} cross-validation). The windows of each person then appear in both the
training and the test set. To show how much this matters, the starting method was evaluated
both ways (Table~\ref{tab:gap}, Fig.~\ref{fig:gap}). Window-wise, it classifies
\R{DSoneRestSegWinAcc}\% of DS-1 rest windows and \R{DStwoSegWinAcc}\% of DS-2 windows
correctly; subject-wise, \R{DSoneRestSubjWinAcc}\% and \R{DStwoSubjWinAcc}\%.

The reason is that the features describe the individual more than the disorder. A
1-nearest-neighbour classifier trained to recognise \emph{which participant} a DS-1 rest window
comes from is correct for \R{DSoneRestProbeAcc}\% of windows (chance \R{DSoneRestProbeChance}\%).
Under window-wise cross-validation the nearest training window of a test window is therefore
almost always a window of the same person, and the classifier copies that person's label.
Window-wise accuracy measures how well new windows of \emph{known} people are recognised, which
is not the task of a detection system. All results in the rest of this report are subject-wise.

\begin{table}[H]
\caption{The starting method (SMVMD features + mRMR + KNN) under window-wise and subject-wise
cross-validation. Window accuracy (\%); brackets: 95\% CI over subjects.}
\label{tab:gap}\centering
\gentable{tab_gap}\\[2pt]
{\footnotesize \dstwonote}
\end{table}

\begin{figure}[H]\centering
\genfig[0.85\textwidth]{segment_vs_subject_gap}
\caption{Window-level (left) and subject-level (right) accuracy of the starting method,
window-wise versus subject-wise. Error bars: 95\% CI over subjects; dashed line: chance.}
\label{fig:gap}
\end{figure}
